\chapter{Software Implementation}
\label{app:software_implementation}

This appendix documents the implementation layer used to generate the numerical results reported in the thesis. Its purpose is not to reproduce the full source code, but to make explicit the computational choices that affect reproducibility: the programming environment, the role of automatic differentiation, the separation between experiment scripts and reusable modules, and the repository in which the implementation was developed.

\section{Computational Environment}

All experiments were implemented in Python. The final neural workflows use PyTorch as the main deep-learning framework, together with NumPy, SciPy, pandas, scikit-learn, and matplotlib for numerical computation, data handling, evaluation, and figure generation. PyTorch is especially useful in this project because the same automatic differentiation engine supports both standard network training and derivative-based diagnostics. In the ANN branch, it is used to train the supervised implied-volatility surrogate and to compute derivatives of the learned surface when Greeks are evaluated. In the PINN branch, it is also used to differentiate the learned price surface with respect to the PDE variables and to build the residual terms that define the physics-informed loss.

The experiments were developed on Apple Silicon hardware with an Apple M4 Max chip, 36 GB of unified memory, and macOS 26.6. PyTorch device selection was handled automatically in the training scripts, using the Apple Metal backend (\texttt{mps}) when available and falling back to CPU otherwise. Numerical reference calculations and benchmark metrics that require double precision were evaluated on CPU when the selected backend did not support the required dtype.

The implementation follows the methodological split used in the thesis. The ANN/CaNN branch is supervised and works mainly in implied-volatility space: synthetic labels are generated through the Heston-COS and Brent pipeline, the ANN surrogate is trained on those labels, and the frozen surrogate is then reused as a fast evaluator in the calibration problem. The PINN branch is implemented separately and works in price space: the baseline model is trained from the Heston PDE residual, terminal condition, and boundary conditions, while the Sobolev and Sobolev+ACV variants add derivative-aware corrections or analytic structure without turning the method into the ANN calibration pipeline.

\section{Code Excerpts}

Only short source-code fragments are included here, because the relevant methodological definitions are already given in Chapters~\ref{ch:ann_surrogate} and~\ref{ch:pinn_greeks}. The objective of these excerpts is to show where the mathematical operations described in the thesis enter the implementation.

\begin{lstlisting}[language=Python, caption={Sobolev derivative loss used in the ANN implied-volatility surrogate. The network is differentiated with respect to its normalized inputs and the resulting gradients are converted back to raw derivative units before comparison with the reference derivative targets.}, label={lst:ann_sobolev_derivative_loss}]
x_req = x_model.detach().clone().requires_grad_(True)
y_model = model(x_req)

grad_model = torch.autograd.grad(
    outputs=y_model,
    inputs=x_req,
    grad_outputs=torch.ones_like(y_model),
    create_graph=True,
    retain_graph=True,
)[0]

idx = torch.as_tensor(
    derivative_indices,
    dtype=torch.long,
    device=x_model.device,
)
x_std_t = torch.as_tensor(x_std, dtype=x_model.dtype, device=x_model.device)
y_std_t = torch.as_tensor(y_std, dtype=x_model.dtype, device=x_model.device)

raw_scale = y_std_t / torch.clamp(x_std_t[idx], min=1.0e-12)
pred_derivatives_raw = (
    grad_model.index_select(dim=1, index=idx)
    * raw_scale.reshape(1, -1)
)

target = target_derivatives_raw.to(dtype=x_model.dtype, device=x_model.device)
scales = torch.clamp(derivative_scales.reshape(1, -1), min=1.0e-12)
loss = torch.mean(
    ((pred_derivatives_raw - target) / scales) ** 2
)
\end{lstlisting}

Listing~\ref{lst:ann_sobolev_derivative_loss} corresponds to the Sobolev term used by the ANN implied-volatility surrogate. The important implementation point is the change of units: if the network receives normalized inputs and, possibly, a normalized output, the derivative returned by automatic differentiation is not yet the physical derivative used in the loss. The factor \(y_{\mathrm{std}}/x_{\mathrm{std},j}\) converts the derivative with respect to the model input back to the raw derivative direction \(j\), after which each derivative component is normalized by its robust scale.

\begin{lstlisting}[language=Python, caption={Construction of the Heston PDE residual used by the baseline PINN. The code shows the moneyness-coordinate version used in the price-space PINN loss, omitting input checks and optional residual normalization for clarity.}, label={lst:pinn_heston_pde_residual}]
x_interior = x_interior.requires_grad_(True)
x_net = apply_input_affine(x_interior, input_affine)
u = model(x_net)

grads = gradient(y=u, x=x_interior)
u_tau = grads[:, 0:1]
u_m = grads[:, 1:2]
u_v = grads[:, 2:3]

grad_u_m = gradient(y=u_m, x=x_interior)
grad_u_v = gradient(y=u_v, x=x_interior)
u_mm = grad_u_m[:, 1:2]
u_mv = grad_u_m[:, 2:3]
u_vv = grad_u_v[:, 2:3]

m = x_interior[:, 1:2]
v = x_interior[:, 2:3]
rho = x_interior[:, 3:4]
kappa = x_interior[:, 4:5]
gamma = x_interior[:, 5:6]
bar_v = x_interior[:, 6:7]
r = x_interior[:, 7:8]

diffusion_m = 0.5 * v * (m**2) * u_mm
mixed = rho * gamma * v * m * u_mv
diffusion_v = 0.5 * (gamma**2) * v * u_vv
drift_m = r * m * u_m
drift_v = kappa * (bar_v - v) * u_v
discount = r * u

residual = (
    u_tau
    - diffusion_m
    - mixed
    - diffusion_v
    - drift_m
    - drift_v
    + discount
)
loss_pde = torch.mean(residual**2)
\end{lstlisting}

Listing~\ref{lst:pinn_heston_pde_residual} shows the corresponding construction for the baseline PINN. The tensor \(x_{\mathrm{interior}}\) contains interior collocation points ordered as \((\tau,m,v,\rho,\kappa,\gamma,\bar v,r)\). The derivatives of the neural price surface are obtained with nested calls to automatic differentiation, and the residual is then assembled term by term according to the Heston PDE written in the moneyness coordinate. In the actual training code, this residual can also be scaled locally before computing the mean-square PDE loss, but the baseline mathematical structure is the one shown above.

\section{Repository}

The code used to generate datasets, train the neural models, run CaNN calibration experiments, evaluate Greeks, and reproduce the figures was maintained in a version-controlled repository during the development of the thesis. The repository contains the reusable modules, experiment scripts, configuration files, trained-run metadata, and plotting utilities associated with the reported benchmarks. The public repository is currently available at \url{https://github.com/jenriquezafra/TFM}.
